\section{Independent word series from a multiplicative shift}\label{sec:words}

We first encode free polynomials by formal power series without losing any
nonzero polynomial. No size bound is needed in this section. Fix a field $F$,
let $t$ be an indeterminate over $F$, and put $K=F(t)$. For a second formal
variable $z$, define
\[
 \sigma:K((z))\longrightarrow K((z)),\qquad
 \sigma(g)(z)=g(tz).
\]
This automorphism fixes $K$ and preserves both $K[[z]]$ and $K(z)$, where
rational functions are embedded by Laurent expansion at $z=0$.
Since $t^r\ne1$ for every nonzero integer $r$, coefficient comparison gives
\begin{equation}\label{eq:fixed-field}
 K((z))^\sigma=K.
\end{equation}
In particular, this remains true when the characteristic divides $r$.

For $1\le i\le n$, set
\begin{equation}\label{eq:letters}
 a_i=(1+t)^i,\qquad b_i(z)=\frac{z}{1-a_i z}.
\end{equation}
Given $g\in K[[z]]$, define $\J_i(g)$ to be the unique series
$h\in zK[[z]]$ satisfying
\begin{equation}\label{eq:shift-antiderivative}
 (\sigma-1)h=b_i g.
\end{equation}
Indeed, the right side has zero constant coefficient, and the coefficient
of $z^r$ in $h$, for $r\ge1$, is obtained by dividing the corresponding coefficient of
$b_i g$ by $t^r-1$. Thus $\J_i$ is $K$-linear and
\begin{equation}\label{eq:order-raising}
 \J_i\bigl(z^jK[[z]]\bigr)\subseteq z^{j+1}K[[z]]\qquad(j\ge0).
\end{equation}
The role of the transcendental parameter is already visible: all these
divisions are valid in every characteristic.

Let $\{1,\ldots,n\}^*$ denote the set of finite words, including the empty
word $\eps$. For $w=i_1\cdots i_\ell$, write
$x_w=x_{i_1}\cdots x_{i_\ell}$, with $x_\eps=1$. Define
\begin{equation}\label{eq:word-series}
 h_\eps=1,\qquad h_{iw}=\J_i(h_w),
\end{equation}
where $iw$ means that $i$ is prefixed to $w$.
Our objective is to prove that these series are linearly independent, even
with rational-function coefficients. The relevant obstruction is that a
rational difference cannot have an isolated pole in an infinite shift orbit.
Such pole-orbit obstructions are central to rational $q$-summation;
see Chen and Singer~\cite[Section~2.3]{ChenSinger2012} for a systematic
residue criterion in characteristic zero. The elementary form needed here
has the following proof in arbitrary characteristic.

\begin{lemma}[A rational difference obstruction]\label{lem:poles}
With $K$, $\sigma$, and $b_1,\ldots,b_n$ as above, suppose
\[
 \sum_{i=1}^n c_i b_i=\sigma(R)-R,
 \qquad c_i\in K,\quad R\in K(z).
\]
Then $c_1=\cdots=c_n=0$.
\end{lemma}

\begin{proof}
The unique finite pole of $b_i$ is $\alpha_i=(1+t)^{-i}$, a nonzero
element of $K$. The orbits
\[
 \mathcal O_i=\{t^j\alpha_i:j\in\mathbb Z\}
\]
are infinite and pairwise disjoint. To check the latter assertion, an orbit
intersection would give $(1+t)^{i-i'}=t^a$ for an integer $a$.
The order at $t=0$ forces $a=0$, and the order at $t=-1$ then forces
$i=i'$. Orders of rational functions are integers; this reasoning is
unchanged when the characteristic divides $i-i'$.

Suppose $c_i\ne0$. The left side has exactly one pole on $\mathcal O_i$.
The rational function $R$ has only finitely many poles there.
If it has none, neither does $\sigma(R)-R$. Otherwise let $j_-$ and $j_+$
be the smallest and largest indices of its poles in this orbit.
Substitution $z\mapsto tz$ moves each pole from index $j$ to index $j-1$.
At index $j_--1$, only $\sigma(R)$ has a pole; at index $j_+$, only $-R$
has one. These are two distinct poles of $\sigma(R)-R$, a contradiction.
The argument does not depend on the pole orders or on cancellation at
interior indices; Figure~\ref{fig:poles} illustrates precisely this point.
\end{proof}

\begin{figure}[tb]
\centering
\begin{tikzpicture}[x=1.1cm,y=0.75cm,font=\small]
 \foreach \y/\lab in {2/{$R$},1/{$\sigma(R)$},0/{$\sigma(R)-R$}} {
  \node[anchor=east] at (-0.55,\y) {\lab};
  \draw[gray!60] (-0.25,\y)--(4.45,\y);
  \foreach \j in {0,...,4} \draw[gray!60] (\j,\y-0.07)--(\j,\y+0.07);
 }
 \foreach \j in {1,3,4} \fill[linkblue] (\j,2) circle (2.5pt);
 \foreach \j in {0,2,3} \fill[linkblue] (\j,1) circle (2.5pt);
 \foreach \j in {0,1,2,4} \fill[linkblue] (\j,0) circle (2.5pt);
 \draw[linkblue,fill=white,thick] (3,0) circle (2.5pt);
 \draw[linkblue,thick] (0,0) circle (5pt);
 \draw[linkblue,thick] (4,0) circle (5pt);
 \draw[-{Stealth[length=2mm]},gray] (1,1.77)--(0,1.23);
 \foreach \j in {0,...,4} \node[below] at (\j,-0.15) {$\j$};
 \node[anchor=west,align=left] at (4.85,0) {interior poles may cancel;\\the two ends cannot};
\end{tikzpicture}
\caption{Pole indices in one orbit $\{t^j\alpha:j\in\mathbb Z\}$,
illustrated with pole set $\{1,3,4\}$ for $R$.
The shift lowers indices by one. The open point allows cancellation at
an overlapping pole, while the circled endpoint poles always survive.
The proof uses the same endpoint argument for every finite nonempty pole set.}
\label{fig:poles}
\end{figure}
\FloatBarrier

The proof below is a shift analogue of the minimal-relation argument for
differential word series in Deneufch\^atel, Duchamp, Hoang Ngoc Minh, and
Solomon~\cite[Theorem~1]{Deneufchatel2011}. We include the argument:
$\sigma-1$ is not a derivation, and the field here may have positive
characteristic.

\begin{proposition}[Word-series independence]\label{thm:words}
For every field $F$ and every positive integer $n$, the series
$(h_w)_{w\in\{1,\ldots,n\}^*}$ in~\eqref{eq:word-series} are linearly
independent over $F(t)(z)$.
\end{proposition}

\begin{proof}
Suppose there is a nontrivial finite relation
\begin{equation}\label{eq:minimal-relation}
 \sum_w R_w h_w=0,\qquad R_w\in K(z).
\end{equation}
Choose one with smallest maximum occurring word length $\ell$, then with
fewest nonzero coefficients at length $\ell$. Since $h_\eps=1$, we have
$\ell\ge1$. Divide by one nonzero length-$\ell$ coefficient to make it $1$.

Apply $\sigma$ to~\eqref{eq:minimal-relation} and subtract the original
relation. By~\eqref{eq:shift-antiderivative}, the coefficient of $h_w$ in
the resulting relation is
\begin{equation}\label{eq:new-relation}
 \widetilde R_w=\sigma(R_w)-R_w+
          \sum_{i=1}^n b_i\sigma(R_{iw}),
\end{equation}
where missing coefficients are zero. This relation has no words longer
than $\ell$, introduces no new length-$\ell$ coefficient, and removes
the normalized coefficient. If any $\widetilde R_w$ were nonzero, one
of the two minimality conditions would fail. Therefore all of them vanish.

For $|w|=\ell$, Equation~\eqref{eq:new-relation} gives
$\sigma(R_w)=R_w$, so~\eqref{eq:fixed-field} implies $R_w\in K$.
For each $|w|=\ell-1$, it then gives
\[
 \sum_{i=1}^n b_i R_{iw}=\sigma(-R_w)-(-R_w).
\]
Lemma~\ref{lem:poles} forces every $R_{iw}$ to vanish. Every word of
length $\ell$ occurs in exactly one such family, contradicting the
choice of the original relation.
\end{proof}

For a free polynomial $f=\sum_w f_w x_w$ over $F$, put
\begin{equation}\label{eq:hf}
 h_f=\sum_w f_w h_w=f(\J_1,\ldots,\J_n)1.
\end{equation}
The rightmost operator acts first, so the second equality follows from
the prefix convention in~\eqref{eq:word-series}. Proposition~\ref{thm:words}
shows that $f\ne0$ implies $h_f\ne0$. To turn this encoding into finite
matrices, we must now bound how far into $h_f$ its first nonzero coefficient
can occur. Formula size supplies that bound through a finite matrix system.
